% --- Archon physics formalization source begin ---
% archon:physics
% archon:covers PhyXMiniProblems/problem_phyx_mini_0266.lean
% archon:source-report reports/phyx_mini/problem_phyx_mini_0266.source.json

\chapter{Physics problem phyx\_mini\_0266}
\label{ch:PhyXMiniProblems_problem_phyx_mini_0266}

\paragraph{Problem source.}
\#\# Physical scenario

In figure, a stick of length \$L = 1.85\textbackslash{} m\$ oscillates as a physical pendulum.

\#\# Question

What is that least period?

\#\# Answer choices

- A: A: 1.4 s

- B: B: 1.8 s

- C: C: 2.1 s

- D: D: 2.4 s

\#\# Figure caption (auxiliary; use the image as primary evidence)

The image shows a diagram of a red rod pivoted at point O, which is attached to a horizontal surface at the top of the figure. The rod is inclined at an angle and appears to rotate around point O, indicated by curved arrows.

The rod is labeled with two specific points:

- Point near the pivot marked as O.
- A point near the center of the rod.

Distance measurements are marked alongside the rod:

- The top portion of the rod, from the pivot to the middle point, is labeled \textbackslash{}(L/2\textbackslash{}).
- The bottom portion, from the middle point to the bottom of the rod, is also labeled \textbackslash{}(L/2\textbackslash{}).

An additional horizontal measurement is labeled \textbackslash{}(x\textbackslash{}), marking the distance from the pivot down to below the midpoint of the rod.

Text and symbols present include:
- The letter "O" near the pivot point.
- "L/2" indicating equal segments along the rod.
- "x" representing a horizontal distance.

The diagram appears to illustrate concepts related to physics, such as rotation or balance.

\#\# Dataset metadata

Resonance and Harmonics; Physical Model Grounding Reasoning, Numerical Reasoning

\paragraph{Recorded answer/context.}
C: C: 2.1 s

\paragraph{Figure/image path.}
/root/proposal\_for\_physic/science-mango/phyx\_mini\_run/phyx\_data/test\_image/266.png

\paragraph{Formalization target.}
create a compiling Lean file with sorry bodies at `PhyXMiniProblems/problem\_phyx\_mini\_0266.lean`. The Lean declarations must preserve the physical quantities, dimensions or dimensional roles, figure labels, governing-law hypotheses, and final relation expressed by this problem.
Use Mathlib/Physlib names found through LeanExplore where available. If a physics API is missing, introduce faithful local abstractions rather than scalar placeholder aliases.

\begin{theorem}[Physics formalization target]
\label{thm:physics:phyx_mini_0266:target}
\lean{PhyXMiniProblems.ProblemPhyXMini0266.leastPeriod_is_answer_C}
\uses{def:physics:phyx-mini-0266:phyxminiproblems-problemphyxmini0266-timeinseconds, def:physics:phyx-mini-0266:phyxminiproblems-problemphyxmini0266-uniformrodphysicalpendulum, def:physics:phyx-mini-0266:phyxminiproblems-problemphyxmini0266-isadmissiblepivotoffsetmeters, def:physics:phyx-mini-0266:phyxminiproblems-problemphyxmini0266-matchesproblemandfiguredata, def:physics:phyx-mini-0266:phyxminiproblems-problemphyxmini0266-satisfiesuniformrodphysicalpendulumlaws, def:physics:phyx-mini-0266:phyxminiproblems-problemphyxmini0266-achievableperiodseconds, def:physics:phyx-mini-0266:phyxminiproblems-problemphyxmini0266-minimizingpivotoffsetmeters, def:physics:phyx-mini-0266:phyxminiproblems-problemphyxmini0266-closedformleastperiodseconds, def:physics:phyx-mini-0266:phyxminiproblems-problemphyxmini0266-answerperiodseconds, def:physics:phyx-mini-0266:phyxminiproblems-problemphyxmini0266-isuniqueclosestanswerchoice}
The assigned autoformalize agent should translate this physics problem into Lean declarations in the covered file, with theorem and lemma proof bodies written as `by sorry`.
\end{theorem}
\begin{proof}
This is an autoformalization task, not a proof task. Produce statements that can later be proved by the physics prover without weakening the physical model.
\end{proof}

% --- Archon declaration topology begin ---
\section{Declaration topology}
\begin{definition}[moment Of Inertia Dimension]
  \label{def:physics:phyx-mini-0266:phyxminiproblems-problemphyxmini0266-momentofinertiadimension}
  \lean{PhyXMiniProblems.ProblemPhyXMini0266.momentOfInertiaDimension}
  The physical dimension of a moment of inertia, mass times length squared
\end{definition}
\begin{proof}
  This is the declared model definition of moment Of Inertia Dimension.
\end{proof}
\begin{definition}[acceleration Dimension]
  \label{def:physics:phyx-mini-0266:phyxminiproblems-problemphyxmini0266-accelerationdimension}
  \lean{PhyXMiniProblems.ProblemPhyXMini0266.accelerationDimension}
  The physical dimension of an acceleration, length divided by time squared
\end{definition}
\begin{proof}
  This is the declared model definition of acceleration Dimension.
\end{proof}
\begin{definition}[Length Quantity]
  \label{def:physics:phyx-mini-0266:phyxminiproblems-problemphyxmini0266-lengthquantity}
  \lean{PhyXMiniProblems.ProblemPhyXMini0266.LengthQuantity}
  A nonnegative physical length, independent of the chosen unit system
\end{definition}
\begin{proof}
  This is the declared model definition of Length Quantity.
\end{proof}
\begin{definition}[Mass Quantity]
  \label{def:physics:phyx-mini-0266:phyxminiproblems-problemphyxmini0266-massquantity}
  \lean{PhyXMiniProblems.ProblemPhyXMini0266.MassQuantity}
  A nonnegative physical mass, independent of the chosen unit system
\end{definition}
\begin{proof}
  This is the declared model definition of Mass Quantity.
\end{proof}
\begin{definition}[Acceleration Quantity]
  \label{def:physics:phyx-mini-0266:phyxminiproblems-problemphyxmini0266-accelerationquantity}
  \lean{PhyXMiniProblems.ProblemPhyXMini0266.AccelerationQuantity}
  \uses{def:physics:phyx-mini-0266:phyxminiproblems-problemphyxmini0266-accelerationdimension}
  A nonnegative physical acceleration magnitude
\end{definition}
\begin{proof}
  This is the declared model definition of Acceleration Quantity.
\end{proof}
\begin{definition}[Moment Of Inertia Quantity]
  \label{def:physics:phyx-mini-0266:phyxminiproblems-problemphyxmini0266-momentofinertiaquantity}
  \lean{PhyXMiniProblems.ProblemPhyXMini0266.MomentOfInertiaQuantity}
  \uses{def:physics:phyx-mini-0266:phyxminiproblems-problemphyxmini0266-momentofinertiadimension}
  A nonnegative physical moment of inertia
\end{definition}
\begin{proof}
  This is the declared model definition of Moment Of Inertia Quantity.
\end{proof}
\begin{definition}[Time Quantity]
  \label{def:physics:phyx-mini-0266:phyxminiproblems-problemphyxmini0266-timequantity}
  \lean{PhyXMiniProblems.ProblemPhyXMini0266.TimeQuantity}
  A nonnegative physical duration
\end{definition}
\begin{proof}
  This is the declared model definition of Time Quantity.
\end{proof}
\begin{definition}[length In Meters]
  \label{def:physics:phyx-mini-0266:phyxminiproblems-problemphyxmini0266-lengthinmeters}
  \lean{PhyXMiniProblems.ProblemPhyXMini0266.lengthInMeters}
  \uses{def:physics:phyx-mini-0266:phyxminiproblems-problemphyxmini0266-lengthquantity}
  Read a physical length in metres
\end{definition}
\begin{proof}
  This is the declared model definition of length In Meters.
\end{proof}
\begin{definition}[mass In Kilograms]
  \label{def:physics:phyx-mini-0266:phyxminiproblems-problemphyxmini0266-massinkilograms}
  \lean{PhyXMiniProblems.ProblemPhyXMini0266.massInKilograms}
  \uses{def:physics:phyx-mini-0266:phyxminiproblems-problemphyxmini0266-massquantity}
  Read a physical mass in kilograms
\end{definition}
\begin{proof}
  This is the declared model definition of mass In Kilograms.
\end{proof}
\begin{definition}[acceleration In Meters Per Second Squared]
  \label{def:physics:phyx-mini-0266:phyxminiproblems-problemphyxmini0266-accelerationinmeterspersecondsquared}
  \lean{PhyXMiniProblems.ProblemPhyXMini0266.accelerationInMetersPerSecondSquared}
  \uses{def:physics:phyx-mini-0266:phyxminiproblems-problemphyxmini0266-accelerationquantity}
  Read an acceleration magnitude in metres per second squared
\end{definition}
\begin{proof}
  This is the declared model definition of acceleration In Meters Per Second Squared.
\end{proof}
\begin{definition}[moment Of Inertia In Kilogram Meter Squared]
  \label{def:physics:phyx-mini-0266:phyxminiproblems-problemphyxmini0266-momentofinertiainkilogrammetersquared}
  \lean{PhyXMiniProblems.ProblemPhyXMini0266.momentOfInertiaInKilogramMeterSquared}
  \uses{def:physics:phyx-mini-0266:phyxminiproblems-problemphyxmini0266-momentofinertiaquantity}
  Read a moment of inertia in kilogram metre squared
\end{definition}
\begin{proof}
  This is the declared model definition of moment Of Inertia In Kilogram Meter Squared.
\end{proof}
\begin{definition}[time In Seconds]
  \label{def:physics:phyx-mini-0266:phyxminiproblems-problemphyxmini0266-timeinseconds}
  \lean{PhyXMiniProblems.ProblemPhyXMini0266.timeInSeconds}
  \uses{def:physics:phyx-mini-0266:phyxminiproblems-problemphyxmini0266-timequantity}
  Read a physical duration in seconds
\end{definition}
\begin{proof}
  This is the declared model definition of time In Seconds.
\end{proof}
\begin{definition}[Rod Figure Point]
  \label{def:physics:phyx-mini-0266:phyxminiproblems-problemphyxmini0266-rodfigurepoint}
  \lean{PhyXMiniProblems.ProblemPhyXMini0266.RodFigurePoint}
  Distinguished points drawn on the rod in the primary figure
\end{definition}
\begin{proof}
  This is the declared model definition of Rod Figure Point.
\end{proof}
\begin{definition}[Rod Figure Text Label]
  \label{def:physics:phyx-mini-0266:phyxminiproblems-problemphyxmini0266-rodfiguretextlabel}
  \lean{PhyXMiniProblems.ProblemPhyXMini0266.RodFigureTextLabel}
  Literal symbolic labels printed in the primary figure
\end{definition}
\begin{proof}
  This is the declared model definition of Rod Figure Text Label.
\end{proof}
\begin{definition}[Uniform Rod Pendulum Figure]
  \label{def:physics:phyx-mini-0266:phyxminiproblems-problemphyxmini0266-uniformrodpendulumfigure}
  \lean{PhyXMiniProblems.ProblemPhyXMini0266.UniformRodPendulumFigure}
  \uses{def:physics:phyx-mini-0266:phyxminiproblems-problemphyxmini0266-lengthquantity, def:physics:phyx-mini-0266:phyxminiproblems-problemphyxmini0266-rodfigurepoint, def:physics:phyx-mini-0266:phyxminiproblems-problemphyxmini0266-rodfiguretextlabel}
  The labelled geometry visible in the physical-pendulum figure
\end{definition}
\begin{proof}
  This is the declared model definition of Uniform Rod Pendulum Figure.
\end{proof}
\begin{definition}[Uniform Rod Physical Pendulum]
  \label{def:physics:phyx-mini-0266:phyxminiproblems-problemphyxmini0266-uniformrodphysicalpendulum}
  \lean{PhyXMiniProblems.ProblemPhyXMini0266.UniformRodPhysicalPendulum}
  \uses{def:physics:phyx-mini-0266:phyxminiproblems-problemphyxmini0266-lengthquantity, def:physics:phyx-mini-0266:phyxminiproblems-problemphyxmini0266-massquantity, def:physics:phyx-mini-0266:phyxminiproblems-problemphyxmini0266-accelerationquantity, def:physics:phyx-mini-0266:phyxminiproblems-problemphyxmini0266-momentofinertiaquantity, def:physics:phyx-mini-0266:phyxminiproblems-problemphyxmini0266-timequantity, def:physics:phyx-mini-0266:phyxminiproblems-problemphyxmini0266-uniformrodpendulumfigure}
  The physical system and its observables. Neither the pivot-dependent moment of inertia nor either period observable is defined by a formula here; the governing laws imposing those formulas are separate hypotheses below. The first input to each function field is the SI readout of the figure coordinate x. The second input to the amplitude-dependent period is the positive, dimensionless angular amplitude measured in radians
\end{definition}
\begin{proof}
  This is the declared model definition of Uniform Rod Physical Pendulum.
\end{proof}
\begin{definition}[Is Admissible Pivot Offset Meters]
  \label{def:physics:phyx-mini-0266:phyxminiproblems-problemphyxmini0266-isadmissiblepivotoffsetmeters}
  \lean{PhyXMiniProblems.ProblemPhyXMini0266.IsAdmissiblePivotOffsetMeters}
  \uses{def:physics:phyx-mini-0266:phyxminiproblems-problemphyxmini0266-lengthinmeters, def:physics:phyx-mini-0266:phyxminiproblems-problemphyxmini0266-uniformrodphysicalpendulum}
  An admissible pivot is strictly away from the center of mass and no farther from it than either end of the rod
\end{definition}
\begin{proof}
  This is the declared model definition of Is Admissible Pivot Offset Meters.
\end{proof}
\begin{definition}[Matches Problem And Figure Data]
  \label{def:physics:phyx-mini-0266:phyxminiproblems-problemphyxmini0266-matchesproblemandfiguredata}
  \lean{PhyXMiniProblems.ProblemPhyXMini0266.MatchesProblemAndFigureData}
  \uses{def:physics:phyx-mini-0266:phyxminiproblems-problemphyxmini0266-lengthinmeters, def:physics:phyx-mini-0266:phyxminiproblems-problemphyxmini0266-accelerationinmeterspersecondsquared, def:physics:phyx-mini-0266:phyxminiproblems-problemphyxmini0266-uniformrodphysicalpendulum}
  Problem and primary-image readouts. The 9.80 m/s² value is the standard near-Earth textbook gravity calibration implicit in the recorded numerical answer. No period or preferred pivot position occurs in these data
\end{definition}
\begin{proof}
  This is the declared model definition of Matches Problem And Figure Data.
\end{proof}
\begin{definition}[Satisfies Uniform Rod Physical Pendulum Laws]
  \label{def:physics:phyx-mini-0266:phyxminiproblems-problemphyxmini0266-satisfiesuniformrodphysicalpendulumlaws}
  \lean{PhyXMiniProblems.ProblemPhyXMini0266.SatisfiesUniformRodPhysicalPendulumLaws}
  \uses{def:physics:phyx-mini-0266:phyxminiproblems-problemphyxmini0266-lengthinmeters, def:physics:phyx-mini-0266:phyxminiproblems-problemphyxmini0266-massinkilograms, def:physics:phyx-mini-0266:phyxminiproblems-problemphyxmini0266-accelerationinmeterspersecondsquared, def:physics:phyx-mini-0266:phyxminiproblems-problemphyxmini0266-momentofinertiainkilogrammetersquared, def:physics:phyx-mini-0266:phyxminiproblems-problemphyxmini0266-timeinseconds, def:physics:phyx-mini-0266:phyxminiproblems-problemphyxmini0266-uniformrodphysicalpendulum, def:physics:phyx-mini-0266:phyxminiproblems-problemphyxmini0266-isadmissiblepivotoffsetmeters}
  The governing physical laws, stated for every admissible figure coordinate x: a uniform thin rod has center moment m L² / 12; the scalar parallel-axis law gives I\_O = I\_cm + m x²; the zero-amplitude limiting period is T₀ = 2 π sqrt (I\_O / (m g x)); the actual amplitude-dependent period approaches T₀ from positive angular amplitudes, with error o(amplitude). The final clause is an explicit local asymptotic contract for the textbook small-angle approximation;
\end{definition}
\begin{proof}
  This is the declared model definition of Satisfies Uniform Rod Physical Pendulum Laws.
\end{proof}
\begin{definition}[Achievable Period Seconds]
  \label{def:physics:phyx-mini-0266:phyxminiproblems-problemphyxmini0266-achievableperiodseconds}
  \lean{PhyXMiniProblems.ProblemPhyXMini0266.AchievablePeriodSeconds}
  \uses{def:physics:phyx-mini-0266:phyxminiproblems-problemphyxmini0266-timeinseconds, def:physics:phyx-mini-0266:phyxminiproblems-problemphyxmini0266-uniformrodphysicalpendulum, def:physics:phyx-mini-0266:phyxminiproblems-problemphyxmini0266-isadmissiblepivotoffsetmeters}
  The set of small-angle periods obtainable by choosing a pivot on the pictured center-to-end half of the rod
\end{definition}
\begin{proof}
  This is the declared model definition of Achievable Period Seconds.
\end{proof}
\begin{definition}[minimizing Pivot Offset Meters]
  \label{def:physics:phyx-mini-0266:phyxminiproblems-problemphyxmini0266-minimizingpivotoffsetmeters}
  \lean{PhyXMiniProblems.ProblemPhyXMini0266.minimizingPivotOffsetMeters}
  \uses{def:physics:phyx-mini-0266:phyxminiproblems-problemphyxmini0266-lengthinmeters, def:physics:phyx-mini-0266:phyxminiproblems-problemphyxmini0266-uniformrodphysicalpendulum}
  The pivot-to-center offset predicted to minimize the period. Its minimizing property remains a conclusion of the theorem below
\end{definition}
\begin{proof}
  This is the declared model definition of minimizing Pivot Offset Meters.
\end{proof}
\begin{definition}[closed Form Least Period Seconds]
  \label{def:physics:phyx-mini-0266:phyxminiproblems-problemphyxmini0266-closedformleastperiodseconds}
  \lean{PhyXMiniProblems.ProblemPhyXMini0266.closedFormLeastPeriodSeconds}
  \uses{def:physics:phyx-mini-0266:phyxminiproblems-problemphyxmini0266-lengthinmeters, def:physics:phyx-mini-0266:phyxminiproblems-problemphyxmini0266-accelerationinmeterspersecondsquared, def:physics:phyx-mini-0266:phyxminiproblems-problemphyxmini0266-uniformrodphysicalpendulum}
  The closed-form candidate for the least period. The theorem below must still prove both its attainability and its lower-bound property
\end{definition}
\begin{proof}
  This is the declared model definition of closed Form Least Period Seconds.
\end{proof}
\begin{definition}[Answer Choice]
  \label{def:physics:phyx-mini-0266:phyxminiproblems-problemphyxmini0266-answerchoice}
  \lean{PhyXMiniProblems.ProblemPhyXMini0266.AnswerChoice}
  The four period readouts printed as answer choices
\end{definition}
\begin{proof}
  This is the declared model definition of Answer Choice.
\end{proof}
\begin{definition}[answer Period Seconds]
  \label{def:physics:phyx-mini-0266:phyxminiproblems-problemphyxmini0266-answerperiodseconds}
  \lean{PhyXMiniProblems.ProblemPhyXMini0266.answerPeriodSeconds}
  \uses{def:physics:phyx-mini-0266:phyxminiproblems-problemphyxmini0266-answerchoice}
  The duration in seconds printed beside each answer choice
\end{definition}
\begin{proof}
  This is the declared model definition of answer Period Seconds.
\end{proof}
\begin{definition}[Is Unique Closest Answer Choice]
  \label{def:physics:phyx-mini-0266:phyxminiproblems-problemphyxmini0266-isuniqueclosestanswerchoice}
  \lean{PhyXMiniProblems.ProblemPhyXMini0266.IsUniqueClosestAnswerChoice}
  \uses{def:physics:phyx-mini-0266:phyxminiproblems-problemphyxmini0266-answerchoice, def:physics:phyx-mini-0266:phyxminiproblems-problemphyxmini0266-answerperiodseconds}
  The selected choice is strictly closer to the computed period than every distinct listed choice
\end{definition}
\begin{proof}
  This is the declared model definition of Is Unique Closest Answer Choice.
\end{proof}
% --- Archon declaration topology end ---
% --- Archon physics formalization source end ---
